\textbf{Phần 1. Trắc nghiệm lựa chọn}
\setcounterpageref{ex}{0}

\begin{ex} %[3D1V02]
Tính nhanh giá trị của biểu thức: $125 + 348 + 275 + 152$.
\choice
{$800$}
{\True $900$}
{$850$}
{$1000$}
\loigiai{
Sử dụng tính chất giao hoán và kết hợp của phép cộng để nhóm các số tròn trăm:\\
$$(125 + 275) + (348 + 152) = 400 + 500 = 900$$
\textbf{Đáp án của câu hỏi là B}
}
\end{ex}

\begin{ex} %[3D1V02]
Trong một phép cộng, nếu số hạng thứ nhất tăng thêm $35$ đơn vị và số hạng thứ hai giảm đi $18$ đơn vị thì tổng mới thay đổi như thế nào so với tổng ban đầu?
\choice
{Giảm đi $17$ đơn vị}
{Tăng thêm $53$ đơn vị}
{Giảm đi $53$ đơn vị}
{\True Tăng thêm $17$ đơn vị}
\loigiai{
Khi số hạng thứ nhất tăng thêm $35$ đơn vị thì tổng tăng $35$ đơn vị.\\
Khi số hạng thứ hai giảm đi $18$ đơn vị thì tổng giảm $18$ đơn vị.\\
Vì $35 > 18$ nên tổng mới tăng thêm số đơn vị là:
$$35 - 18 = 17\text{ (đơn vị)}$$
\textbf{Đáp án của câu hỏi là D}
}
\end{ex}

\begin{ex} %[3D1V02]
Trong một phép trừ có hiệu bằng $142$. Biết tổng của số bị trừ, số trừ và hiệu bằng $560$. Số bị trừ của phép trừ đó là
\choice
{\True $280$}
{$140$}
{$418$}
{$210$}
\loigiai{
Ta có hệ thức: $\text{Số bị trừ} = \text{Số trừ} + \text{Hiệu}$.\\
Do đó:\\
$$\text{Số bị trừ} + \text{Số trừ} + \text{Hiệu} = \text{Số bị trừ} + \text{Số bị trừ} = 2 \times \text{Số bị trừ}$$
Số bị trừ là:
$$560 : 2 = 280$$
\textbf{Đáp án của câu hỏi là A}
}
\end{ex}

\begin{ex} %[3D1V02]
Cho phép cộng đặt tính: $\overline{4a6} + \overline{25b} = 743$. Giá trị của biểu thức $a + b$ là
\choice
{$13$}
{$14$}
{\True $15$}
{$16$}
\loigiai{
Xét phép cộng theo từng cột từ phải sang trái:\\
- Hàng đơn vị: $6 + b$ có tận cùng là $3$, do $6 + 7 = 13$ nên $b = 7$ (nhớ $1$ sang hàng chục).\\
- Hàng chục: $a + 5 + 1 = 14$ nên $a + 6 = 14 \Rightarrow a = 8$ (nhớ $1$ sang hàng trăm).\\
- Hàng trăm: $4 + 2 + 1 = 7$ (thỏa mãn).\\
Vậy $a = 8$ và $b = 7$. Ta có: $a + b = 8 + 7 = 15$.\\
\textbf{Đáp án của câu hỏi là C}
}
\end{ex}

\begin{ex} %[3D1V02]
Một phép trừ có hiệu bằng $240$. Nếu số bị trừ tăng thêm $40$ đơn vị và số trừ tăng thêm $15$ đơn vị thì hiệu mới là
\choice
{$295$}
{\True $265$}
{$215$}
{$185$}
\loigiai{
Khi số bị trừ tăng thêm $40$ đơn vị thì hiệu tăng $40$ đơn vị.\\
Khi số trừ tăng thêm $15$ đơn vị thì hiệu giảm đi $15$ đơn vị.\\
Hiệu mới là:
$$240 + 40 - 15 = 265$$
\textbf{Đáp án của câu hỏi là B}
}
\end{ex}

\begin{ex} %[3D1V02]
Hiệu giữa số lớn nhất có 3 chữ số khác nhau và số bé nhất có 3 chữ số khác nhau là
\choice
{$880$}
{$875$}
{$890$}
{\True $885$}
\loigiai{
- Số lớn nhất có 3 chữ số khác nhau là: $987$.\\
- Số bé nhất có 3 chữ số khác nhau là: $102$.\\
Hiệu giữa hai số là:
$$987 - 102 = 885$$
\textbf{Đáp án của câu hỏi là D}
}
\end{ex}

\begin{ex} %[3D1V02]
Thư viện trường có ba tủ sách. Tủ thứ nhất có $245$ cuốn sách. Tủ thứ hai có nhiều hơn tủ thứ nhất $35$ cuốn sách nhưng lại ít hơn tủ thứ ba $20$ cuốn sách. Cả ba tủ sách có tất cả số cuốn sách là
\choice
{\True $825$ cuốn sách}
{$805$ cuốn sách}
{$845$ cuốn sách}
{$795$ cuốn sách}
\loigiai{
Số sách ở tủ thứ hai là:
$$245 + 35 = 280\text{ (cuốn)}$$
Số sách ở tủ thứ ba là:
$$280 + 20 = 300\text{ (cuốn)}$$
Cả ba tủ có tất cả số cuốn sách là:
$$245 + 280 + 300 = 825\text{ (cuốn)}$$
\textbf{Đáp án của câu hỏi là A}
}
\end{ex}

\begin{ex} %[3D1V02]
Tính giá trị của biểu thức: $P = 782 - (245 + 182)$.
\choice
{$345$}
{$365$}
{\True $355$}
{$335$}
\loigiai{
Bỏ dấu ngoặc hoặc trừ liên tiếp để tính nhanh:\\
$$P = 782 - 182 - 245 = 600 - 245 = 355$$
\textbf{Đáp án của câu hỏi là C}
}
\end{ex}

\begin{ex} %[3D1V02]
Hai số có tổng bằng $640$. Biết số bé là một số có hai chữ số và khi viết thêm chữ số $5$ vào bên trái số bé ta được số lớn. Số lớn là
\choice
{$550$}
{\True $570$}
{$560$}
{$580$}
\loigiai{
Khi viết thêm chữ số $5$ vào bên trái một số có hai chữ số thì số đó tăng thêm $500$ đơn vị.\\
Do đó, số lớn hơn số bé $500$ đơn vị.\\
Số bé là:
$$(640 - 500) : 2 = 70$$
Số lớn là:
$$70 + 500 = 570$$
\textbf{Đáp án của câu hỏi là B}
}
\end{ex}

\begin{ex} %[3D1V02]
Tìm chữ số $x$ thỏa mãn phép trừ: $\overline{7x4} - 289 = 465$.
\choice
{$3$}
{$4$}
{$6$}
{\True $5$}
\loigiai{
Ta có:
$$\overline{7x4} = 465 + 289 = 754$$
So sánh hai vế, ta được chữ số $x = 5$.\\
\textbf{Đáp án của câu hỏi là D}
}
\end{ex}

\begin{ex} %[3D1V02]
Một cửa hàng ngày thứ nhất bán được $318\text{ m}$ vải. Ngày thứ hai bán được ít hơn ngày thứ nhất $45\text{ m}$ vải. Ngày thứ ba bán được ít hơn ngày thứ hai $20\text{ m}$ vải. Cả ba ngày cửa hàng bán được số mét vải là
\choice
{\True $844\text{ m}$}
{$824\text{ m}$}
{$864\text{ m}$}
{$834\text{ m}$}
\loigiai{
Ngày thứ hai bán được số mét vải là:
$$318 - 45 = 273\text{ (m)}$$
Ngày thứ ba bán được số mét vải là:
$$273 - 20 = 253\text{ (m)}$$
Cả ba ngày bán được số mét vải là:
$$318 + 273 + 253 = 844\text{ (m)}$$
\textbf{Đáp án của câu hỏi là A}
}
\end{ex}

\begin{ex} %[3D1V02]
Từ bốn chữ số $1; 3; 6; 8$, lập số lớn nhất có 3 chữ số khác nhau và số bé nhất có 3 chữ số khác nhau. Hiệu của hai số lập được là
\choice
{$717$}
{$737$}
{\True $727$}
{$729$}
\loigiai{
- Số lớn nhất có 3 chữ số khác nhau từ 4 chữ số trên là: $863$.\\
- Số bé nhất có 3 chữ số khác nhau từ 4 chữ số trên là: $136$.\\
Hiệu của hai số là:
$$863 - 136 = 727$$
\textbf{Đáp án của câu hỏi là C}
}
\end{ex}

\textbf{Phần 2. Trắc nghiệm đúng sai}
\setcounterpageref{ex}{0}

\begin{ex} %[3D1V02]
Xét tính đúng sai của các nhận định sau về phép cộng và phép trừ trong phạm vi 1000:
\choiceTF
{\True Khi thêm vào số hạng này bao nhiêu đơn vị và bớt ở số hạng kia bấy nhiêu đơn vị thì tổng không thay đổi.}
{Trong một phép trừ, số bị trừ luôn luôn lớn hơn số trừ.}
{\True Nếu $a + b = 500$ thì $(a + 50) + (b + 50) = 600$.}
{\True Hiệu giữa số chẵn lớn nhất có 3 chữ số và số chẵn bé nhất có 3 chữ số bằng $898$.}
\loigiai{
\begin{itemize}
    \item Ý a) Đúng: Ví dụ $(a + k) + (b - k) = a + b + k - k = a + b$.
    \item Ý b) Sai: Số bị trừ có thể bằng số trừ (khi đó hiệu bằng 0, ví dụ $500 - 500 = 0$).
    \item Ý c) Đúng: Cả hai số hạng đều tăng 50 đơn vị nên tổng tăng: $50 + 50 = 100$. Tổng mới là $500 + 100 = 600$.
    \item Ý d) Đúng: Số chẵn lớn nhất có 3 chữ số là 998, số chẵn bé nhất có 3 chữ số là 100. Hiệu là $998 - 100 = 898$.
\end{itemize}
\begin{itemize}
    \item a) \textbf{Đúng}.
    \item b) \textbf{Sai}.
    \item c) \textbf{Đúng}.
    \item d) \textbf{Đúng}.
\end{itemize}
}
\end{ex}

\begin{ex} %[3D1V02]
Cho ba số $A = 345, B = 456, C = 199$. Xét tính đúng sai của các khẳng định sau:
\choiceTF
{\True Tổng $A + B + C$ là một số tròn nghìn.}
{Hiệu $(B - A)$ lớn hơn hiệu $(A - C)$.}
{\True Nếu bớt ở cả ba số mỗi số $10$ đơn vị thì tổng mới giảm đi $30$ đơn vị so với tổng ban đầu.}
{Số $B$ lớn hơn tổng của hai số $A$ và $C$.}
\loigiai{
\begin{itemize}
    \item Ý a) Đúng: $A + B + C = 345 + 456 + 199 = 1000$ (là số tròn nghìn).
    \item Ý b) Sai: $B - A = 456 - 345 = 111$; $A - C = 345 - 199 = 146$. Vì $111 < 146$ nên $B - A < A - C$.
    \item Ý c) Đúng: Cả ba số cùng giảm 10 đơn vị nên tổng giảm $10 \times 3 = 30$ đơn vị.
    \item Ý d) Sai: $A + C = 345 + 199 = 544$. Vì $456 < 544$ nên $B$ bé hơn tổng của $A$ và $C$.
\end{itemize}
\begin{itemize}
    \item a) \textbf{Đúng}.
    \item b) \textbf{Sai}.
    \item c) \textbf{Đúng}.
    \item d) \textbf{Đúng}.
\end{itemize}
}
\end{ex}

\textbf{Phần 3. Trả lời ngắn}
\setcounterpageref{ex}{0}

\begin{ex} %[3D1V02]
Tính nhanh giá trị của biểu thức: $234 + 187 + 166 + 213$.
\shortans{800}
\loigiai{
\textbf{Đáp số:} $800$\\
Nhóm các số tròn trăm:\\
$$(234 + 166) + (187 + 213) = 400 + 400 = 800$$
}
\end{ex}

\begin{ex} %[3D1V02]
Tìm một số tự nhiên, biết rằng nếu lấy số đó cộng với $125$ rồi trừ đi $85$ thì được kết quả là $400$.
\shortans{360}
\loigiai{
\textbf{Đáp số:} $360$\\
Gọi số cần tìm là $x$. Theo bài ra ta có:\\
$$x + 125 - 85 = 400$$
$$x + 40 = 400$$
$$x = 400 - 40 = 360$$
}
\end{ex}

\begin{ex} %[3D1V02]
Hai thùng chứa tất cả $450\text{ lít}$ dầu. Nếu người ta chuyển $30\text{ lít}$ dầu từ thùng thứ nhất sang thùng thứ hai thì lúc này tổng số dầu ở cả hai thùng là bao nhiêu lít?
\shortans{450}
\loigiai{
\textbf{Đáp số:} $450$\\
Khi chuyển dầu qua lại giữa hai thùng thì tổng số lít dầu ở cả hai thùng không thay đổi, vẫn là $450\text{ lít}$.
}
\end{ex}

\begin{ex} %[3D1V02]
Một cuộn dây điện dài $240\text{ m}$. Lần thứ nhất người thợ cắt đi $65\text{ m}$ dây. Lần thứ hai người đó cắt nhiều hơn lần thứ nhất $15\text{ m}$ dây. Hỏi sau hai lần cắt, cuộn dây còn lại bao nhiêu mét?
\shortans{95}
\loigiai{
\textbf{Đáp số:} $95$\\
Lần thứ hai cắt số mét dây điện là:
$$65 + 15 = 80\text{ (m)}$$
Cả hai lần cắt tất cả số mét dây điện là:
$$65 + 80 = 145\text{ (m)}$$
Sau hai lần cắt, cuộn dây còn lại số mét là:
$$240 - 145 = 95\text{ (m)}$$
}
\end{ex}

\begin{ex} %[3D1V02]
Tìm một số có 3 chữ số có chữ số hàng trăm là $3$. Biết rằng nếu xóa bỏ chữ số $3$ này đi thì số đó giảm đi bao nhiêu đơn vị?
\shortans{300}
\loigiai{
\textbf{Đáp số:} $300$\\
Số có dạng $\overline{3ab} = 300 + \overline{ab}$.\\
Khi xóa chữ số 3 ở hàng trăm thì được số mới là $\overline{ab}$.\\
Số đó giảm đi: $(300 + \overline{ab}) - \overline{ab} = 300$ đơn vị.
}
\end{ex}

\begin{ex} %[3D1V02]
Trong một phép trừ, biết tổng của số bị trừ, số trừ và hiệu bằng $720$. Tìm số bị trừ của phép trừ đó.
\shortans{360}
\loigiai{
\textbf{Đáp số:} $360$\\
Vì $\text{Số bị trừ} = \text{Số trừ} + \text{Hiệu}$ nên:\\
$$\text{Số bị trừ} + \text{Số trừ} + \text{Hiệu} = 2 \times \text{Số bị trừ} = 720$$
Số bị trừ là:
$$720 : 2 = 360$$
}
\end{ex}
